\section{Bölüm sonu çözümlü problemler}

\begin{enumerate}
  \item \textbf{Standartlaştırılmış alfa.} Altı maddelik bir ölçekte ortalama
  maddeler arası korelasyon $\bar r=0{,}30$ bulunmuştur. Standartlaştırılmış
  alfayı hesaplayın.

  \textbf{Çözüm:}
  \[
  \alpha_{\mathrm{std}}
  =\frac{6(0{,}30)}{1+5(0{,}30)}
  =\frac{1{,}80}{2{,}50}=0{,}72.
  \]
  Bu değer tek başına ölçeğin tek boyutlu veya geçerli olduğunu göstermez.

  \item \textbf{Ters madde hatası.} Beş maddelik bir ölçeğin alfası 0,31'dir.
  Dördüncü maddenin diğer maddelerle korelasyonları negatiftir ve madde
  olumsuz yönde yazılmıştır. Araştırmacının izlemesi gereken sırayı yazın.

  \textbf{Çözüm:} Önce özgün puanlama anahtarı incelenir; yalnız madde metnine
  bakılarak karar verilmez. Madde gerçekten tersse doğru aralık formülüyle
  yeniden kodlanır. Değer aralığı ve frekanslar kontrol edilir; ardından alfa
  ve düzeltilmiş madde--toplam korelasyonları yeniden hesaplanır. İlk katsayı
  ölçeğin zayıflığı olarak değil, olası puanlama hatasının işareti olarak ele
  alınır.

  \item \textbf{Madde silme kararı.} Toplam alfa 0,82, bir madde silindiğinde
  alfa 0,83'tür. Maddenin düzeltilmiş madde--toplam korelasyonu 0,38'dir.
  Madde silinmeli midir?

  \textbf{Çözüm:} Yalnız 0,01'lik alfa artışı silme kararı için yeterli
  değildir. Maddenin ölçek kapsamındaki özgün katkısı, ifadesinin açıklığı,
  alt boyut yapısı, başka örneklemlerdeki davranışı ve silmenin içerik
  geçerliğine etkisi incelenmelidir. Madde kuramsal olarak önemliyse korunması
  daha uygun olabilir.

  \item \textbf{Uygun kanıtı seçme.} Üç danışman aynı görüşme kayıtlarını
  bağımsız olarak 0--4 arasında puanlamaktadır. Araştırma ekibi yalnız
  Cronbach alfa raporlamıştır. Sorunu açıklayın.

  \textbf{Çözüm:} Temel hata kaynağı madde örneklemi değil değerlendiricidir.
  Sürekli/sıralı puanın kullanımına ve değerlendirici tasarımına uygun bir ICC
  modeli veya uygun uyum katsayısı seçilmelidir. Model türü, tek değerlendirici
  ya da ortalama puan ve mutlak uyum/tutarlılık hedefi açıkça belirtilmelidir.
\end{enumerate}